% =====================================================================
% Draft model section for Section 3 (after "The audited system").
% Every fact is traced to pva-experiments @ ecbd49a in model-section-facts.md.
% Lines starting with "% CHECK" are for the authors and must be resolved or deleted.
% CHECK: add \label{eq:repair} to the COR'/CAIL repair equation in S5 of paper.tex (it has no
% label; it is no longer Eq. (2) once this section precedes it). The paper is also being
% edited (new epsilon-multiplicity equations), so refer to equations by label only.
% =====================================================================
\subsection{The model}\label{sec:model}
This subsection states the audited model as the reference implementation computes it. The
membership functions and rule bases are those of~\cite{rabiei2026resilient}; the full rule
tables are given there.

\subsubsection{Decision and inputs}
Let $J=\{1,\dots,m\}$ index the affected people and $I=\{1,\dots,n\}$ the relief centres, with
$0<n_{\mathrm{dir}}<m$. A decision $d=(D,c)$ is a set $D\subset J$ of exactly $n_{\mathrm{dir}}$
directed people and an assignment $c\colon D\to I$ of each directed person to one centre;
equivalently $x_{j,i}\in\{0,1\}$ with $\sum_{i}x_{j,i}\le 1$ for every $j$ and
$\sum_{j,i}x_{j,i}=n_{\mathrm{dir}}$. People outside $D$ are not directed. The published model
has no capacity constraint: any number of people may be sent to one centre.

Table~\ref{tab:inputs} lists the inputs. Ordinal fields enter the model through fixed scores.
Two scores are derived. The age score is $\mathrm{AS}_j=(25-a_j)/25$ for $a_j<25$, $0$ for
$25\le a_j\le 60$, $(a_j-60)/30$ for $60<a_j\le 90$ and $1$ above; the vulnerability score and
the roadworthiness score are
\begin{equation}\label{eq:vs-rws}
\mathrm{VS}_j=\tfrac14\bigl(\mathrm{AS}_j+\mathrm{DS}_j+\mathrm{IL}_j+\mathrm{LS}_j\bigr),
\qquad
\mathrm{RWS}_{j,i}=1-\tfrac12\bigl(\mathrm{RCS}_{j,i}+\mathrm{PHS}_{j,i}\bigr),
\end{equation}
the weighted means of~\cite{rabiei2026resilient} with all weights equal to one. $\mathrm{RWS}$
is written so that $1$ is a clear, safe route, which is the reading its rule base requires.

\begin{table}[H]
\centering
\caption{Model inputs, their encoding and the instance generator. Scores are listed in the
order of the levels. All draws are independent; $\mathrm{U}$ is uniform, $\mathrm{B}$ the beta
distribution, and probabilities follow the levels in order.\label{tab:inputs}}
\small
\begin{tabularx}{\textwidth}{@{}l>{\raggedright\arraybackslash}X>{\raggedright\arraybackslash}X>{\raggedright\arraybackslash}X@{}}
\toprule
Symbol & Input (unit) & Domain or levels $\to$ score & Generator\\
\midrule
$a_j$ & age (years) & $[0,120]$ & $\mathrm{U}(5,90)$, 1 decimal\\
$\mathrm{DS}_j$ & disability status & none, minor, severe $\to 0, 0.5, 1$ & $0.70, 0.20, 0.10$\\
$\mathrm{IL}_j$ & injury level & none, minor, moderate, serious, life-threatening $\to 0, 0.25, 0.5, 0.75, 1$ & $0.40, 0.25, 0.20, 0.10, 0.05$\\
$\mathrm{LS}_j$ & living status & with support, alone $\to 0, 1$ & $0.65, 0.35$\\
$\mathrm{IDL}_j$ & infrastructure damage at the person's location (\%) & $[0,100]$ & $100\cdot\mathrm{B}(2.5,2)$\\
$\mathrm{RTR}_j$ & resource time remaining (h) & $[0,48]$ & $\mathrm{U}(1,48)$\\
$\mathrm{COR}_i$ & centre occupancy rate (\% of capacity in use) & $[0,100]$ & $\mathrm{U}(20,90)$\\
$\mathrm{RDR}_i$ & resource depletion rate (\% of resources per hour) & $[0,100]$ & $\mathrm{U}(10,80)$\\
$\mathrm{TD}_{j,i}$ & travel duration from person $j$ to centre $i$ (min) & $[0,180]$ & $\min\{180,\max\{2,\,100\cdot\mathrm{B}(1.5,4)\}\}$\\
$\mathrm{RCS}_{j,i}$ & road condition & clear, partially blocked, blocked $\to 0.1, 0.5, 1$ & $0.45, 0.40, 0.15$\\
$\mathrm{PHS}_{j,i}$ & possible hazard & none, minor, moderate, significant, extreme $\to 0.05, 0.25, 0.5, 0.75, 1$ & $0.30, 0.30, 0.25, 0.10, 0.05$\\
\bottomrule
\end{tabularx}
\end{table}

\subsubsection{Fuzzy inference systems}
Three Mamdani systems produce the scores of the three-objective model
(Table~\ref{tab:fis}). Every input has three terms on its universe: with breakpoints
$p_0<\dots<p_4$, a left-shoulder trapezoid $(p_0,p_0,p_1,p_2)$, a triangle $(p_1,p_2,p_3)$ and a
right-shoulder trapezoid $(p_2,p_3,p_4,p_4)$. Every output has five terms on $[0,100]$: very low
$(0,0,10,30)$, low $(10,30,50)$, moderate $(30,50,70)$, high $(50,70,90)$ and very high
$(70,90,100,100)$. The rule bases are complete grids, one rule per combination of input terms.
Rules fire with the minimum of their antecedent memberships, clip their consequent term at
that strength, and are aggregated by the maximum; the output is the centroid of the
aggregate. Universes are sampled at 101 points (travel duration at 181), crisp inputs are
clipped to the interior of their universe before inference, and a system that fires no rule
returns the mid-point $50$.

\begin{table}[H]
\centering
\caption{The fuzzy inference systems. Breakpoints $(p_0,\dots,p_4)$ as defined in the text;
all outputs on $[0,100]$ with five terms.\label{tab:fis}}
\small
\begin{tabularx}{\textwidth}{@{}l>{\raggedright\arraybackslash}Xcl@{}}
\toprule
System & Inputs: universe, terms, breakpoints & Rules & Output\\
\midrule
$\mathrm{FIS}_1$ & $\mathrm{VS}$: $[0,1]$, low/medium/high, $(0,.25,.5,.75,1)$;
$\mathrm{IDL}$: $[0,100]$, minor/moderate/severe, $(0,25,50,75,100)$;
$\mathrm{RTR}$: $[0,48]$\,h, short/moderate/long, $(0,12,24,36,48)$ & 27 & $\mathrm{ULPP}_j$\\
$\mathrm{FIS}_2$ & $\mathrm{TD}$: $[0,180]$\,min, short/moderate/long, $(0,30,60,90,180)$;
$\mathrm{RWS}$: $[0,1]$, low/moderate/high, $(0,.25,.5,.75,1)$ & 9 & $\mathrm{TIL}_{j,i}$\\
$\mathrm{FIS}_3$ & $\mathrm{COR}$: $[0,100]$, low/moderate/high, $(0,25,50,75,100)$;
$\mathrm{RDR}$: $[0,100]$, slow/moderate/rapid, $(0,25,50,75,100)$;
$\mathrm{TD}$ as in $\mathrm{FIS}_2$ & 27 & $\mathrm{CAIL}_{j,i}$\\
\midrule
$\mathrm{FIS}_{2a}$ & $\mathrm{RCS}$: $[0,1]$, $(0,.2,.5,.8,1)$; $\mathrm{PHS}$: $[0,1]$, low
$(0,0,.2,.4)$, moderate $(.2,.5,.8)$, high $(.6,.8,1,1)$ & 9 & $\mathrm{TRD}_{j,i}$\\
$\mathrm{FIS}_{2b}$ & $\mathrm{TD}$ as in $\mathrm{FIS}_2$ & 3 & $\mathrm{RPD}_{j,i}$\\
\bottomrule
\end{tabularx}
\end{table}

Two rules of $\mathrm{FIS}_1$ show the direction of its scale: if $\mathrm{VS}$ is high and
$\mathrm{IDL}$ is severe and $\mathrm{RTR}$ is short, $\mathrm{ULPP}$ is very high; if
$\mathrm{VS}$ is low and $\mathrm{IDL}$ is minor and $\mathrm{RTR}$ is long, $\mathrm{ULPP}$ is
very low. In $\mathrm{FIS}_3$, low occupancy, slow depletion and a short trip give very low
$\mathrm{CAIL}$, and high occupancy with rapid depletion gives very high $\mathrm{CAIL}$ at
every travel duration.
% CHECK (author decision): ULPP_j rises with VS_j and IDL_j and falls with RTR_j, and f_1 below
% is minimised over the directed people. The model therefore prefers to direct people with
% LOW ULPP. On the small instance the exact equal-weight optimum directs people whose mean crisp
% priority is 0.31 (all people: 0.39; crisp rule: 0.53), and 32 of the 40 units by which the
% crisp rule trails the optimum come from f_1 (72.1 vs 40.0). See facts file, item F-ULPP.
% Decide whether the paper states this, before reviewers infer it from this section.

\subsubsection{Objectives}
With $\mathrm{FIS}_1$--$\mathrm{FIS}_3$ the per-person and per-pair scores are
\begin{equation}\label{eq:scores}
\mathrm{ULPP}_j=\mathrm{FIS}_1(\mathrm{VS}_j,\mathrm{IDL}_j,\mathrm{RTR}_j),\quad
\mathrm{TIL}_{j,i}=\mathrm{FIS}_2(\mathrm{TD}_{j,i},\mathrm{RWS}_{j,i}),\quad
\mathrm{CAIL}_{j,i}=\mathrm{FIS}_3(\mathrm{COR}_i,\mathrm{RDR}_i,\mathrm{TD}_{j,i}),
\end{equation}
and the three objectives, all minimised, are plain means over the directed people:
\begin{equation}\label{eq:objectives}
f_1(d)=\frac{1}{n_{\mathrm{dir}}}\sum_{j\in D}\mathrm{ULPP}_j,\qquad
f_2(d)=\frac{1}{n_{\mathrm{dir}}}\sum_{j\in D}\mathrm{TIL}_{j,c(j)},\qquad
f_3(d)=\frac{1}{n_{\mathrm{dir}}}\sum_{j\in D}\mathrm{CAIL}_{j,c(j)}.
\end{equation}
In the published model $\mathrm{CAIL}_{j,i}$ is the per-pair $\mathrm{FIS}_3$ score at the
centre's input occupancy $\mathrm{COR}_i$, which no decision changes. In the repaired model of
S5 the same rule base is evaluated at the post-allocation occupancy
$\mathrm{COR}'_i(\ell_i)$ of Eq.~(\ref{eq:repair}), where $\ell_i=|\{j\in D: c(j)=i\}|$ and
$\mathrm{cap}_i=\lceil\kappa\,n_{\mathrm{dir}}/n\rceil$ persons; nothing else changes. The
\emph{quality} of a decision is $Q(d)=f_1(d)+f_2(d)+f_3(d)$, always computed on clean inputs.
In the four-objective variant of~\cite{rabiei2026resilient}, $f_2$ is replaced by the means of
$\mathrm{TRD}_{j,c(j)}=\mathrm{FIS}_{2a}(\mathrm{RCS},\mathrm{PHS})$ and
$\mathrm{RPD}_{j,c(j)}=\mathrm{FIS}_{2b}(\mathrm{TD})$ (Table~\ref{tab:fis}). All scores are
computed once per instance and cached, so evaluating a decision costs $O(n_{\mathrm{dir}})$
look-ups.

\subsubsection{Encoding and operators}
A chromosome has $2n_{\mathrm{dir}}$ genes. Genes $1,\dots,n_{\mathrm{dir}}$ are distinct
person indices (a partial permutation of $J$); gene $n_{\mathrm{dir}}+k$ is a real
$r_k\in[0,1)$ that assigns the $k$-th listed person to centre
$1+\min\{\lfloor n\,r_k\rfloor,\,n-1\}$. Any chromosome therefore decodes to a decision with
exactly $n_{\mathrm{dir}}$ directed people, each at one centre, and the order of the pairs
carries no meaning. The initial population draws the person slice as a uniform sample
without replacement and each $r_k$ from $\mathrm{U}[0,1)$.

The three solvers share one $(\mu+\lambda)$ loop with $\lambda=\mu$ and differ only in
survivor selection. Each generation copies the population; there is no separate mating
selection. Consecutive copies are paired, and each pair is crossed with probability $0.7$.
On the person slice each child takes every position from one parent or the other with
probability $0.5$; repeated persons are then dropped (the first occurrence is kept), the list
closes up, and it is refilled at the end with persons drawn at random from those not yet
listed. On the centre slice a blend crossover with $\alpha=0.5$ draws, per gene,
$\gamma\sim\mathrm{U}[-0.5,1.5]$ and sets the children to $(1-\gamma)a+\gamma b$ and
$\gamma a+(1-\gamma)b$, clipped to $[0,1)$. Each copy is then mutated with probability $0.2$.
Mutation acts per gene with probability $0.05$: a person gene is replaced by a person not
currently directed, and the person it displaces becomes the next replacement, so one
mutation changes the directed set by at most one person; a centre gene receives Gaussian
noise with $\sigma=0.1$, clipped to $[0,1)$. The $\mu$ survivors are selected from parents
and offspring together: NSGA-II uses non-dominated sorting with crowding distance; NRGA, as implemented,
fills the next population front by front and samples the last admitted front uniformly at
random, without crowding distance; NSGA-III uses Das--Dennis reference points with $p=4$
divisions ($15$ points for three objectives, $35$ for four)~\cite{deb2014nsga3}. A run stops
after a fixed number of generations and returns the first non-dominated front of its final population, sorted
lexicographically by objective vector. Table~\ref{tab:params} lists the settings of every run.

\begin{table}[H]
\centering
\caption{Solver settings. All runs use the three-objective model unless stated. Seeds are
$1000+7919k$, $k=0,1,\dots$; ``evaluations'' counts $\mu\times$ generations.\label{tab:params}}
\small
\begin{tabularx}{\textwidth}{@{}>{\raggedright\arraybackslash}p{3.6cm}>{\raggedright\arraybackslash}X@{}}
\toprule
Setting & Value\\
\midrule
Population $\mu=\lambda$ & $100$ in all runs, except the equal-evaluation controls ($200$ and $400$)\\
Crossover & probability $0.7$ per pair; person slice uniform ($0.5$) with repair; centre slice blend, $\alpha=0.5$\\
Mutation & probability $0.2$ per individual; per gene $0.05$; person gene: swap-in of an undirected person; centre gene: $\mathcal{N}(0,0.1^2)$, clipped to $[0,1)$\\
Survivor selection & NSGA-II: rank and crowding distance; NRGA: rank, random truncation of the last front; NSGA-III: reference points, $p=4$\\
Termination & fixed number of generations (below)\\
\midrule
Operational budget (S3, S4) & NSGA-II, $150$ generations; $20$ seeds for the seed floor; $8$ (small) and $6$ (large) seeds in the turbulence runs\\
Budget curve (S3) & NSGA-II, $150, 300, 600, 1200, 2400, 4800, 9600$ generations; $5$ seeds above $150$\\
Solver controls (S2, S3) & NRGA and NSGA-III, small instance: $150$ generations ($10$ seeds) and $2400$ generations ($5$ seeds)\\
Population controls (S2, S3) & NSGA-II, small, $240{,}000$ evaluations: $\mu=200$ for $1200$ and $\mu=400$ for $600$ generations; $5$ seeds\\
Instance controls (S2, S3) & NSGA-II, $150$ generations, two further instances per size, $10$ seeds each\\
Envelope (S5) & NSGA-II on the repaired model ($\kappa=2$), seed $1000$; generations $\max\{20,\lfloor T/t_{\mathrm{gen}}\rfloor\}$, where $t_{\mathrm{gen}}$ is the mean time per generation of a 20-generation run and $T=60$ or $300$\,s\\
% CHECK: this row describes run_scaling.py (the r1 envelope). If the revision switches to
% run_envelope.py, replace it by: "NSGA-II, plain and seeded with the fallback and 21 exact
% weight-sweep optima; stopped by a wall-clock deadline T = 60 or 300 s (generation bound 1e9)",
% and restore "or at a wall-clock deadline (S5)" in the Termination row and the text above.
Published experiments~\cite{rabiei2026resilient} & NSGA-II, NRGA, NSGA-III; three and four objectives; $200$ generations; $30$ seeds ($k=0,\dots,29$)\\
\bottomrule
\end{tabularx}
\end{table}

\subsubsection{Reference deciders}
The \emph{exact} decider is the separable optimum of S1 at equal weights: person $j$ costs
$\mathrm{ULPP}_j+\min_i(\mathrm{TIL}_{j,i}+\mathrm{CAIL}_{j,i})$ at its cheapest centre, and the
$n_{\mathrm{dir}}$ cheapest people are directed there; ties go to the lower index. The
\emph{crisp priority} rule uses no fuzzy inference. It ranks people by
\begin{equation}\label{eq:crisp-priority}
\pi_j=\tfrac16\Bigl(\mathrm{AS}_j+\mathrm{DS}_j+\mathrm{IL}_j+\mathrm{LS}_j+\frac{\mathrm{IDL}_j}{100}
+1-\frac{\min\{\mathrm{RTR}_j,72\}}{72}\Bigr),
\end{equation}
directs the $n_{\mathrm{dir}}$ people with the highest $\pi_j$ (ties to the lower index), and
sends each to the centre of lowest crisp cost
\begin{equation}\label{eq:crisp-cost}
k_{j,i}=\tfrac15\Bigl(\frac{\mathrm{TD}_{j,i}}{180}+\frac{\mathrm{COR}_i}{100}+\frac{\mathrm{RDR}_i}{100}
+\mathrm{RCS}_{j,i}+\mathrm{PHS}_{j,i}\Bigr),
\end{equation}
again with ties to the lower index and without regard to load. Both rules are deterministic;
their decisions are scored by the same objectives (\ref{eq:objectives}) as the evolutionary
system's.

\subsubsection{Committal rule and stability measures}
From a front $F$ the committal rule takes
\begin{equation}\label{eq:committal}
d^\ast=\arg\min_{d\in F}\sum_{k=1}^{3}\frac{f_k(d)-\min_{F}f_k}{\max_{F}f_k-\min_{F}f_k},
\end{equation}
where an objective that is constant on $F$ contributes zero, and ties go to the first point
in the sorted front. The alternative rule of S3 takes $\arg\min_{d\in F}\lVert f(d)\rVert_2$
in raw objectives. For a reference decision $A=(D_A,c_A)$ and a comparison decision
$B=(D_B,c_B)$,
\begin{equation}\label{eq:churn}
\chi(A,B)=\frac{\bigl|\{j\in D_A:\ j\notin D_B\ \text{or}\ c_B(j)\neq c_A(j)\}\bigr|}{|D_A|},
\qquad
\mathrm{Jac}(A,B)=\frac{|D_A\cap D_B|}{|D_A\cup D_B|},
\qquad
\Delta Q(A,B)=Q(B)-Q(A).
\end{equation}
The denominator of the churn $\chi$ is the reference decision's directed set: the clean
decision in S4, the lower-indexed seed in S3. A person directed only in $B$ does not enter the
numerator. Because both decisions direct exactly $n_{\mathrm{dir}}$ people, every person in
$D_A\setminus D_B$ is matched by one in $D_B\setminus D_A$, so
$\chi(A,B)=1-|\{j\in D_A\cap D_B: c_A(j)=c_B(j)\}|/n_{\mathrm{dir}}$ and the measure is
symmetric in value. All three measures are computed on clean data.

The seed floor is the mean over the $190$ pairs $P$ of the $20$ clean seed decisions $A_s$,
$\bar\chi_0=|P|^{-1}\sum_{(a,b)\in P,\,a<b}\chi(A_a,A_b)$, with a 95\% interval from $4000$
bootstrap resamples of seeds in which self-pairs are dropped. In S4, realisation $r$ of a
perturbation at level $\lambda$ is drawn from a generator seeded with $1000+r$, the same for
every level, and the system decides on the perturbed instance with each of its $S$ solver
seeds; $A_s$ is the clean decision and $B_{s,r}^\lambda$ the perturbed one for seed $s$. With
$R$ realisations, the churn above the floor and the quality harm are
\begin{equation}\label{eq:above-floor}
\Delta\chi(\lambda)=\frac1R\sum_{r=1}^{R}\frac1S\sum_{s=1}^{S}\chi\bigl(A_s,B_{s,r}^{\lambda}\bigr)-\bar\chi_0,
\qquad
L_r(\lambda)=\frac1S\sum_{s=1}^{S}\Delta Q\bigl(A_s,B_{s,r}^{\lambda}\bigr),
\end{equation}
with $R=12$, $S=8$ (small) and $R=8$, $S=6$ (large). The harm is paired by seed, so its null
is zero and no floor is subtracted from it; the seed floor enters as its scale,
$\bar L(\lambda)/\sigma_Q$, where $\bar L(\lambda)$ is the mean of $L_r(\lambda)$ over
realisations and $\sigma_Q$ the sample standard deviation of $Q(A_s)$ over the $20$ clean seeds.
Intervals for $\bar L$ and for the mean churn bootstrap the realisations ($4000$ resamples).
The trend test is Spearman's $\rho$ between $\lambda\in\{0.05,0.1,0.2,0.4\}$ and $L_r(\lambda)$
over all $(r,\lambda)$, one-sided for $\rho>0$, Holm-adjusted across cells within a system
and size. For the exact decider ($S=1$, $R=50$) the same $\Delta Q$ is its regret.

\subsubsection{Instances and perturbations}
The reference instances are \emph{small} ($n=5$, $m=150$, $n_{\mathrm{dir}}=50$),
\emph{medium} ($8$, $225$, $75$) and \emph{large} ($10$, $300$, $100$), each drawn from the
distributions of Table~\ref{tab:inputs} by one generator per size with base seed $42$, in the
order people, centres, routes. They are fixed files, archived with the code and identified by
their SHA-256 in every manifest.
% CHECK (facts file, item F-SALT): the committed instances were generated with a size salt
% taken from Python's per-process str hash(), so they cannot be regenerated; the working tree
% now salts with CRC-32 (uncommitted), which yields different files. Do not claim the reference
% instances regenerate from the seed. The same applies to the alternate instances (seed43/44)
% and to the r1 envelope instances. Fields are drawn independently: no input
depends on another, and travel durations are drawn per pair, not from a spatial layout, so
they lie in $[2,100]$\,min and the upper clip at $180$ never binds. The only couplings are
structural: $\mathrm{COR}_i$ and $\mathrm{RDR}_i$ are shared by every route into centre $i$,
and $\mathrm{VS}_j$ and $\mathrm{RWS}_{j,i}$ combine several fields by
Eq.~(\ref{eq:vs-rws}). The envelope instances of S5 come from the same generator with
$m\in\{150,300,600,1000,1500,3000\}$, $n\in\{5,10,20,25,40,50\}$,
$n_{\mathrm{dir}}=\lfloor m/3\rfloor$ and base seeds $7+s$. Turbulence (S4) scales noise
and clips perturbed values to fixed ranges: $[0,100]$ for $\mathrm{IDL}$, $\mathrm{COR}$ and
$\mathrm{RDR}$, $[0,72]$\,h for $\mathrm{RTR}$ (values above $48$ enter $\mathrm{FIS}_1$ at the
top of its universe), $[0,180]$\,min for $\mathrm{TD}$ and $[0,120]$ for age. A flip replaces
a categorical level by one of the other levels with equal probability; missing values are
replaced by the instance median (continuous) or mode (categorical).
